% Fragmento público generado desde propietarios íntegros.
% Residencia: 52.10.2.
% Fuente: ../incoming/radion_monografia_20260827/manuscrito/sections/11_modularidad_orbitas.tex:85-187
\subsubsection{Areal Conservation in Kepler and Sommerfeld Orbits}

\paragraph{Areal Equivalence of the Orbital and Oscillatory Realizations}

Kepler's ellipse lives in configuration space and responds to a central
potential. The radion ellipse lives in phase space and responds to an action
form. They are different objects. They share an areal law because both
realizations possess a conserved two-form.

For a Keplerian orbit,
\begin{equation}
r(\phi)=\frac{p}{1+e\cos(\phi-\phi_0)},
\qquad
\frac{\diff A}{\diff t}=\frac{\ell}{2m}.
\label{eq:kepler}
\end{equation}
The third law takes the form
\[
T^2=\frac{4\pi^2\mu}{K}a^3.
\]
For the phase oscillator,
\begin{equation}
J=\oint p\,\diff q=\frac{2\pi E}{\omega}.
\label{eq:osc-action}
\end{equation}
If the elementary cell is \(J=\hfull\), then
\begin{equation}
E=\hret\omega,\qquad \frac{J}{2\pi}=\hret.
\label{eq:E-hbarw}
\end{equation}
This is the clean convergence between Planck action, angular frequency, and radion.

Sommerfeld quantizes action integrals over cycles. The HMT reading adds
the memory of which cycle, sheet, and return survive. On a Möbius band,
one visible loop may carry \(\hfull/2\) and two loops restore \(\hfull\).
That topological factor is not added to the Maslov index of EBK quantization;
they belong to corrections with distinct owners.

\paragraph{Orbital Holonomy and Advanced and Retarded Propagations}

An orbit that accumulates holonomy satisfies a condition of the type
\begin{equation}
\epsilon_\gamma\exp\left(\frac{iJ_\gamma}{\hret}\right)=1.
\label{eq:holonomy-action}
\end{equation}
The sign \(\epsilon_\gamma\) preserves the sheet; \(J_\gamma\) preserves the action. Apsidal advance or retardation is interpreted as a decoupling between visible angular closure and complete return. The radion supplies the oriented unit with which that difference is measured.

\subsubsection{Composition of Schrödinger amplitudes and path summation}

\paragraph{Four realizations of the same phase}

Once \(\hret\) is fixed, distinct conventional equations recognize
aspects of the same transport:
\begin{enumerate}
\item Schrödinger publishes the phase through
\(i\hret\partial_t\psi=H\psi\).
\item Pauli adds the double spin leaf and its magnetic coupling.
\item Dirac linearizes the energy--momentum relation and organizes particle and antiparticle within a causal representation.
\item The path sum assigns to each history the phase
\(\exp(iS[\gamma]/\hret)\).
\end{enumerate}
These theories do not retrospectively generate \(\hret\). They receive the unit of
action and realize different contents of the enriched state.

\paragraph{Phase functional over the path space}

In Feynman's formulation,
\begin{equation}
\mathcal K(b,a)=\int_{\gamma:a\to b}
\exp\left(\frac{i}{\hret}S[\gamma]\right)\mathcal D\gamma.
\label{eq:path-integral}
\end{equation}
An action increment \(\hret\) changes the phase by one radian. An increment
\(\hfull\) changes it by one revolution. The theorem
\(\mathcal A(\theta)=\hret\theta\) converts this relation into area geometry:
the phase of a path can be read as swept action in units of the
radion.

The inversion \(C_M=-\operatorname{rev}\) on dodecaphase words and the identity \eqref{eq:time-reversal} organize conjugate routes. In an absorber realization, the following are distinguished
\[
G_{\mathrm{sym}}=\tfrac12(G_{\mathrm{ret}}+G_{\mathrm{adv}}),
\qquad
G_{\mathrm{odd}}=G_{\mathrm{ret}}-G_{\mathrm{adv}}.
\]
The first sector is even under retarded--advanced interchange; the second is
odd. This reproduces the separation between even energy/form and odd
oriented action of the radion when the transducer preserves the boundary and the
symplectic form.

\begin{statusbox}
Classical absorber theory and the Feynman--Stueckelberg interpretation of
antiparticles are distinct historical realizations. The former uses a
temporally symmetric interaction; the latter reorganizes propagation in
relativistic theory. HMT supplies a common antecedent of reversal, but does not
merge them without a map of propagators.
\end{statusbox}
